\section{Quotient centers and early projection in the free ten-positive
chart}\label{quotient-centers-and-early-projection-in-the-free-ten-positive-chart}

\subsection{Result, domain, and
attribution}\label{result-domain-and-attribution}

This package gives three ordinary polynomials in exactly ten
independently quantified strictly positive integer variables:

\begin{longtable}[]{@{}lrrlr@{}}
\toprule\noalign{}
Graph prefix & Operations & Exact degree & Transport & Root scale \\
\midrule\noalign{}
\endhead
\bottomrule\noalign{}
\endlastfoot
quotient\_t & 97 & 202 & bV & 2 \\
product2 & 111 & 198 & bV & 2 \\
product4 & 119 & 190 & b\^{}2 & 4 \\
\end{longtable}

Every graph is supplied at two separate complete native/refutation
calibrations: historical indexed943/radix67 and current low-padding2963.
The latter has 69,685 raw coefficient bytes, not the frozen4177 tuple's
98,219 bytes. Each of its four integers, native proof, exact942-selector
bijection, physical marker and Pclp controller changes as one coherent
unit. No coefficient port, old word or controller certificate is
transferred alone.

The free chart Y=s, lambda=q, inverse-Lucas correction and unsquared
trace are credited to source172, whose full retained parent is
\texttt{parent/}. Source178 supplies the A-quotient idea and the
full-product Omega4 evaluator; its original proof and algebra are
retained in \texttt{inherited/}. Those incoming scaled-chart profiles
96/220 and116/208 do not improve the prior free-chart 90/206. Our
contributions are the compatible free-chart composition, the cheaper
t=A-1 quotient, first-product truncation, and early p=3 projector
exclusion which newly licenses square transport here. The last step
avoids illicitly importing the scaled-chart q\textbar Y argument.

All binary +,-,* gates, squares and scalar products are charged. Fixed
integer leaves and reuse are free. These are measured tradeoffs, not a
global optimality or literature-priority claim. Ordinary
soundness/completeness, composed with the identified retained syntactic
calibration, proves in ZFC:

\begin{verbatim}
The equation has no positive-integer zero iff Con(ZFC).
\end{verbatim}

No consistency decision, actual zero, unsolvability assertion, new
compiled formal-kernel replay or new encoded internal-ZFC proof object
is claimed. Finite algebra and byte checks do not prove the unbounded
lemmas by testing.

\subsection{1. The complete ordinary
polynomials}\label{the-complete-ordinary-polynomials}

The independent variables in graph order are

\begin{verbatim}
J,F,alpha,i,T,w,s,n,eta,zeta > 0.
\end{verbatim}

Fix the entire chosen tuple beta,K0,MC,MF. Both actual tuples satisfy
beta\textgreater=73 odd, K0\textgreater0 and
0\textless MC\textless beta\textless MF\textless2beta. The strengthened
beta bound is used explicitly in section3; no theorem for arbitrary
beta\textgreater=31 is asserted. MF is the source mask with beta added
exactly once, not the native mask. Define these polynomial
abbreviations, none a new quantified variable:

\begin{verbatim}
h=beta J, q=h+1, C=q-F-alpha, X=qw, Y=s, Ybar=qY,
k=eta+zeta, c=kY+eta,
A=Ybar(X+1)+1, Delta=A^2-1, H=4A-5,
v=X Ybar^2, ell=vk, E=1+ell(ell+k),
R=[h(q-F)+alpha](q^2-1)+(MC+qMF)J,
B=(K0+w)C-F, U=ic^2, a=Delta U^2E, F0=1+a,
G=EF0, f=a+F0, Q=4Delta aF0,
b=cT-R, V=fb-c, K=X(3c-2X)+2,
e6=6(k-R-1)-ell(k^2-4), D=6q^7, M=hDH.
\end{verbatim}

Choose a row below; the formulas are division-free integer polynomials.
For quotient\_t put

\begin{verbatim}
Qt=(kYbar)[ell+k-(kYbar)Ybar],
C2=U(2AE-Qt), Omega=C2 V, r=2, L=bV.
\end{verbatim}

For product2 put

\begin{verbatim}
S=(kYbar)[ell+k-(kYbar)(Ybar+1)], theta=U^2b,
Omega=U{2A[EV-theta(E^2+S^2)]
              +S[4theta E-(b-c)]}, r=2, L=bV.
\end{verbatim}

For product4 put

\begin{verbatim}
S=(kYbar)[ell+k-(kYbar)(Ybar+1)], theta=U^2b,
e2=(kYbar)^2, Rc=theta(3E+e2)-(b-c),
Sc=4E^2+3S^2+e2^2,
Omega=U[2S Rc+A(4EV-theta Sc)], r=4, L=b^2.
\end{verbatim}

In every case set

\begin{verbatim}
Z0=hD[H(n-LB)+K]-H[DB+h(e6-6q^4s)],
Z=Z0-M Omega,
P=Q(Z^2-r^2M^2GV^2)+(r^2M^2G-Z^2).                 (1)
\end{verbatim}

The single equation is P=0. No rational constant, division, recurrence,
radical, floor, variable exponent, hidden witness or oracle is a graph
gate. The quotient\_t implementation shares V in n-V(bB+C2) and computes
2A=(A-1)+(A+1) using two already-live registers.

For proof only put t0=B/h, kappa=K/H, rho=(e6/6)/q\^{}7, sigma=Y/q\^{}3,
W=1+hL and z=Z/M. Before integrality is established these are rational
quantities. Identically,

\begin{verbatim}
z=n+kappa+sigma-Wt0-rho-Omega,
P=(Q-1)Z^2-r^2M^2G(QV^2-1).                    (2)
\end{verbatim}

\subsection{2. Root extraction before sign or denominator
assumptions}\label{root-extraction-before-sign-or-denominator-assumptions}

At every positive input q\textgreater=74, X\textgreater=q,
Y\textgreater=1, c\textgreater=3, A\textgreater q\^{}2, H\textgreater0,
U\textgreater=c\^{}2, f\textgreater2c and Q\textgreater c\^{}2. The
exact identities

\begin{verbatim}
f^2-1=4aF0, Q=Delta(f^2-1)=4Delta^2U^2G,
gcd(E,F0)=1, f-1=2Delta i^2c^4E
\end{verbatim}

are unchanged from the retained parent. V cannot vanish. If P=0, the
integer equation(2), valuations at primes dividing G, and the square
factor r\^{}2 show that G is a square. Thus E=tau\^{}2, F0=f0\^{}2 and
g=tau f0\textgreater0. The parent's rational adjacent-Pell descent,
unchanged by the constant square scale, gives y=\textbar z\textbar/(rg)
an integer and

\begin{verbatim}
QV^2-(Q-1)y^2=1, z=epsilon rgy, epsilon=+/-1.
\end{verbatim}

In particular z is integral. The signed all-index folding theorem
applies with I=2Ug, f\^{}2-Delta I\^{}2=1, Q=(Delta I)\^{}2 and
c\^{}2\textbar I. It gives c=psi\_A(p) for p positive odd;
c\textgreater=3 excludes p=1. Hence

\begin{verbatim}
p>=3, c>=c3=4Delta+3>4q^4, K>0.
\end{verbatim}

All of this is independent of the signs of R,B,C,b,V. No q\textbar Y,
authentication p=R, auxiliary-index-three exclusion or large-rank ratio
bound has been used.

The first square gives k=2psi\_(1+2v)(m), m\textgreater=1. m=1 gives
k=2, eta=zeta=1 and c=2Y+1 strictly between the first two main
ordinates, so m\textgreater=2. The first base 1+2v\textgreater A, and
c\textgreater k. Strict base/index monotonicity therefore gives
m\textless p: if m\textgreater=p,
psi\_(1+2v)(m)\textgreater psi\_A(p)=c, impossible. This m\textless p
fact does NOT use p=R or any mask bound.

\subsection{3. New early projector exclusion of main index
three}\label{new-early-projector-exclusion-of-main-index-three}

Integrality of z already gives H\textbar Z. Reducing the literal
numerator(1) modulo H gives H\textbar hDK. Since D=6q\^{}7 and
gcd(H,q)=1,

\begin{verbatim}
H | 6hK.                                      (3)
\end{verbatim}

The shift M Omega and either transport are multiples of H; they do not
change this step. We retain the entire factor6h.

Suppose p=3. Then m=2 by section2. Write u=X+2-2Y. Using c=4A\^{}2-1,
k=8Xq\^{}2Y\^{}2+4, eta=c-kY gives the exact polynomial identity

\begin{verbatim}
eta=4q^2Y^2(Xu+1)+8qY(X+1)-4Y+3.              (4)
\end{verbatim}

If u\textless=-1, the first term is at most -4q\^{}2Y\^{}2(X-1). Since
qY\textgreater=32, X\textgreater=32, we have
qY(X-1)\textgreater2(X+1)+1, so(4) is negative. If u\textgreater=2,
subtracting k from(4) leaves at least
4q\^{}2Y\^{}2+8qY(X+1)-4Y-1\textgreater0. Both alternatives contradict
0\textless eta\textless k. Thus u is0 or1, so 2Y=X+2 or X+1 and

\begin{verbatim}
H=4qY(X+1)-1>2qX^2.                           (5)
\end{verbatim}

For p=3 the projector is the direct identity

\begin{verbatim}
4K+(X-8)(8X+1)=3XH(4A+5).                   (6)
\end{verbatim}

Together(3),(6) show H divides the positive integer N=6(q-1)(X-8)(8X+1).
Because X\textgreater=32, 0\textless N\textless48qX\^{}2. By(5), j=N/H
is an integer with1\textless=j\textless=23. Modulo q, H congruent -1, X
congruent 0 and N congruent 48, so j congruent -48 modq.
But49\textless=j+48\textless=71\textless q, impossible. Therefore
p\textgreater=5 and c\textgreater c3. This is a universal proof at the
explicitly verified beta\textgreater=73 boundary; finite tests of small
inputs are not its justification.

Now b cannot vanish by the adjacent-square gap argument from the parent.
It has the sign of V. The auxiliary Pell index is odd and at least3. At
index3 its two signs would give respectively

\begin{verbatim}
f(b-4Delta f)=c-c3,
f(b+4Delta f)=c+c3.
\end{verbatim}

Both right sides lie strictly between0 and2c\textless f, impossible.
Thus that index is at least5, and

\begin{verbatim}
|V|>=16Q^2-20Q+5>8Q^2,
|b|>8f^3-1/2>7f^3, f>6g,
y<3|V|/2, |V|<3f|b|/2.                     (7)
\end{verbatim}

This finally licenses square transport. The proof did not silently
exclude an index-three branch before its required early projector
argument. For the two bV rows the weaker inherited bound
\textbar b\textbar\textgreater4Delta f-1 would already suffice:
Delta\textgreater=4 and f\textgreater6g imply 4Delta
f-1\textgreater96g-1\textgreater72g for integer g\textgreater=1. their
validity does not depend on square transport.

\subsection{4. Centers are small, with both signs
covered}\label{centers-are-small-with-both-signs-covered}

Let C0=AUE. Since G=E+Delta U\^{}2E\^{}2,

\begin{verbatim}
0<C0<3g/2.                                  (8)
\end{verbatim}

For quotient\_t set Rt=1+k\^{}2Ybar\^{}2(Ybar\^{}2-1)\textgreater0.
Exact polynomial division by t=A-1 gives E=(A-1)Qt+Rt. Also
Qt=k\^{}2Ybar{[}(X-1)Ybar\^{}2+1{]}\textgreater0, hence
Qt\textless E/(A-1). Consequently 0\textless C2=U(2AE-Qt)\textless2AUE.
Thus Omega has the sign of V and

\begin{verbatim}
2gy+|Omega|<3g|V|+3g|V|=6g|V|<L/2.         (9)
\end{verbatim}

For the product rows, temporarily regard A as an indeterminate T over
Z{[}Ybar,k,U,b,c{]}. Then

\begin{verbatim}
E(T)=e2 T^2+e1 T+e0,
e2=k^2Ybar^2,
e1=k^2Ybar(1-2Ybar-2Ybar^2),
e0=1+k^2Ybar(Ybar+1)[Ybar(Ybar+1)-1],
V(T)=(1+2(T^2-1)U^2E(T))b-c, F(T)=E(T)V(T).
\end{verbatim}

Let Qj be the value at A of the polynomial quotient of F(T) by T\^{}j.
Then exactly

\begin{verbatim}
Omega2=U(2AF-Q1),
Omega4=U(4AF-2Q1-Q3/2).                     (10)
\end{verbatim}

Q3 is coefficientwise even, because the only odd-coefficient candidate
(b-c)E has T-degree at most2. Equations(10) are proof-only descriptions;
the literal evaluators in section1 use no quotient or division gate.

For a polynomial, superscript\# denotes absolute-coefficient evaluation
at A. Put delta=(Ybar+1)/A\textless1/16. The coefficient bounds
\textbar e1\textbar\textless2e2(Ybar+1),
0\textless e0\textless=e2(Ybar+1)\^{}2 give E\#(A)\textless2E. Since
c+1\textless2U\^{}2E and
\textbar V\textbar\textgreater(f-c)\textbar b\textbar\textgreater Delta
U\^{}2E\textbar b\textbar,

\begin{verbatim}
F#(A)<2E(|b|+c)+8(A^2+1)U^2E^2|b|,
F#(A)/|EV|<8+20/Delta<10.
\end{verbatim}

Deletion of low coefficients gives
\textbar Qj\textbar\textless=F\#(A)/A\^{}j. As A\textgreater=16,

\begin{verbatim}
|Omega2-2AUEV|<AUE|V|,
|Omega4-4AUEV|<AUE|V|.
\end{verbatim}

Therefore both have the sign of V, and

\begin{verbatim}
|Omega2|<3AUE|V|,
3AUE|V|<|Omega4|<5AUE|V|.                   (11)
\end{verbatim}

For product2 with L=bV, (8),(11) and \textbar b\textbar\textgreater72g
yield
2gy+\textbar Omega2\textbar\textless(15/2)g\textbar V\textbar\textless L/2.
For product4 with L=b\^{}2, (7),(8),(11) give

\begin{verbatim}
4gy+|Omega4|<(9/4)f|V|
  <(27/8)f^2|b|<b^2/4.                       (12)
\end{verbatim}

The last inequality follows from \textbar b\textbar\textgreater7f\^{}3
and f\textgreater2. All bounds precede sign recovery and cover negative
b,V. We do not call Omega positive until V itself has been proved
positive.

\subsection{5. Signed bootstrap, authentication, and projection
order}\label{signed-bootstrap-authentication-and-projection-order}

The inverse correction satisfies e*=e6/6\textless=k-R-1\textless b for
every integer R, so rho\textless\textbar b\textbar/q\^{}7. In both
weights L\textgreater0 and L\textgreater=\textbar b\textbar. If
B\textless=-1, Wt0=B/h+LB\textless-L. Solving(2) for positive n and
using(9)--(12), positive kappa,sigma, and q\^{}7\textgreater2 gives

\begin{verbatim}
n< L/2-L+|b|/q^7<0,
\end{verbatim}

a contradiction. Thus B\textgreater=0, and (K0+w)C=B+F\textgreater0
gives C\textgreater0. Only now use the selected whole tuple's rank
bound3q\textless R\textless q\^{}4-2. As
c\textgreater4q\^{}4\textgreater2R, b,V are positive. The unchanged
normalized main-index and signed auxiliary-index theorems applied to V
congruent -c modf and V congruent -R modc identify p=R. Mask parity now
gives q even,J odd.

Every Omega is an integer polynomial. Its numerator contribution M Omega
is divisible by D and H, so all subsequent modular reductions are
exactly those of parent/PROOF.md sections4--5. We state their
noncircular order. First k=2psi\_(1+2v)(m) is even and e* is integral.
Literal cancellation of6 in the rational carrier gives
q\^{}4\textbar e*. The full Lucas identity is

\begin{verbatim}
k=2m+8v Cm+32v^2Zm, Cm=binom(m+1,3), Zm>=0,
3[e*-(2m-R-1)]=4v^2[24Zm-Am Sm],
Am=Cm+4vZm, Sm=k^2+k(2m)+(2m)^2-4.
\end{verbatim}

With v=q\^{}3ws\^{}2 the latter gives q\^{}4\textbar e*-(2m-R-1) at
every prime, including3. Since m\textless R and R\textless q\^{}4,
m=(R+1)/2. One does not yet infer q\^{}7\textbar e*. Write R=2r0+1;
r0\textgreater=48. The lambda-uniform ratio theorem at
lambda=q,Y\textgreater=1 gives Y\textgreater X\^{}r0/3 and
X\^{}(r0+1)\textless H. The full-product projector bound, retaining6h,
gives X=2\^{}R. Hence q is a power of2; rank growth supplies
q\^{}6\textbar X and v is divisible by q\^{}8. Only now does the cubic
identity imply q\^{}7\textbar e*. Thus rho and kappa are integral.

The remaining carrier relation yields q\^{}3\textbar Y and h\textbar B,
since W congruent 1 modh. Therefore t0 is a nonnegative integer. The
retained fixed-minus parity, lambda-uniform half-binomial rounding and
fresh native reconstruction at A\_source=Y(X+1)+2 produce actual source
words for the selected whole calibration. This base is not confused with
A=qY(X+1)+1. The native interpretation yields a finite syntactic ZFC
refutation, not merely a conclusion from elementary scalar boxes.

\subsection{6. Constructive strictly positive
completeness}\label{constructive-strictly-positive-completeness}

Conversely, the chosen whole native/refutation calibration supplies
fresh canonical data q=2\^{}a, X=2\^{}R, R congruent 3
mod4,3q+1\textless=R\textless q\^{}4,q\^{}3\textbar Y,C\textgreater0 and
t0=B/h\textgreater=0 integral. Parent section6 is uniform in these data
and constructs the free-chart inputs w=X/q,s=Y, positive eta,zeta,i,T,
with c=psi\_A(R), k=2psi\_(1+2v)((R+1)/2),
g\textgreater0,V\textgreater0,y\textgreater0 and

\begin{verbatim}
g>kappa+sigma+|rho|.
\end{verbatim}

No old packed word or strong period is reused across calibrations. Every
center is positive when V\textgreater0 by section4. Set, at the selected
transport and root scale,

\begin{verbatim}
n=rgy+Omega+(1+hL)t0+rho-kappa-sigma.        (13)
\end{verbatim}

All terms are integers and the preceding strict bound makes
n\textgreater0. The literal numerator satisfies z=rgy, so(1) vanishes.
This construction also handles possibly negative rho. For bV one may
equivalently take n\_new=n\_parent+(r-1)gy+Omega. For b\^{}2 we use(13),
not a claim that subtracting bVB from an arbitrary positive parent n
preserves positivity. All ten independent coordinates have their
declared strict positive domain.

\subsection{7. Entire low-padding2963 native
calibration}\label{entire-low-padding2963-native-calibration}

The new native proof, decoder, language and every historical ordinary
interface it names are retained in native2963/, with267 exact source
pins. Its complete proof is research/low-padding-independent/PROOF.md.
The admissible selector bijection g(0)=0,g(j)=j+1
(1\textless=j\textless=941) removes exactly old forbidden selector1. It
preserves and reflects the full language, including the
complete436-character seed,430 actions, renamed initializer928 and
calibrated final tile941; it is not an unproved canonical subfamily.

There are3297 coefficient slots,976 weighted positions and978 physical
positions. Physical dummy977 and marker1005 have zero coefficient but
remain Boolean positions. A separate kernel monomial copies Start0
to1005. All916 padding tests use bit2 of fields1..916. The sole clause
is1956; marker parity, Boolean temporal count, alignment, carries and
reset tests are proved before the resulting native interpretation.
Soundness needs no controller premise. Completeness uses the same
tuple's discharged Pclp:

\begin{verbatim}
rho_native=2^47, L_native=2963, d=139261,
fixed prime factors47,2963, orders23,2962,
spacing68126, auxiliary prime1362521,
H residues29,2938,492699.
\end{verbatim}

The finite checker reconstructs the exact four unsigned big-endian
integers, not just their boxes; MF\_source adds beta once. Every
raw/compressed hash, bit length, minimal encoding, marker monomial and
physical support is checked. The complete native section12 constructs
fresh canonical outer data matching section6 exactly. The free-chart
endpoint in section5 constructs fresh original source witnesses, so the
native sections5--8 apply. This establishes both interfaces rather than
substituting a new tuple into an old calibration.

The historical radix67 specialization is separately supplied with the
full parent finite recipe and little-endian payloads,295388 raw bytes.
The new tuple is a different specialized polynomial and is not claimed
identical to either radix67 or any older4177/4027 specialization.

\subsection{8. Exact algebra, degree, and complete highest
forms}\label{exact-algebra-degree-and-complete-highest-forms}

The checker independently expands the center over Z{[}A,Ybar,k,U,b,c{]},
using weights(4,2,1,5,4,2). It compares the division-free actual graph
against(10) and the t-quotient identity. Put
Def=Omega\^{}2-r\^{}2GV\^{}2. The weighted degrees of Omega and Def are

\begin{verbatim}
quotient_t:59,106; product2:59,102; product4:59,94.
\end{verbatim}

For quotient\_t the exact root defect is

\begin{verbatim}
C2^2-4G=U^2[4E(Rt-Qt)+Qt^2]-4E.
\end{verbatim}

No positivity of Rt-Qt is required. The full three defect leaders are

\begin{verbatim}
16A^10 U^6 Ybar^10 b^2 k^8-16A^11 U^6 Ybar^8 b^2 k^8,
4A^10 U^6 Ybar^8 b^2 k^8,
8A^8 U^6 Ybar^8 b^2 k^8-48A^6 U^6 Ybar^12 b^2 k^8.
\end{verbatim}

The complete output identity is

\begin{verbatim}
P=Q[Z0^2-2M Z0 Omega+M^2Def]+r^2M^2G-(Z0-M Omega)^2.
\end{verbatim}

Independent degrees are deg(Q,M,G,V)=(72,12,46,36). For bV, degZ0=54;
for b\^{}2, degZ0=22. The corresponding summand bounds
are(180,197,202,70,142),(180,197,198,70,142), and (116,165,190,70,142).
In each case the defect summand uniquely dominates. Let q0=beta
J,k0=eta+zeta. The complete collected top forms are

\begin{verbatim}
[P_quotient_t]_202 = 36864 q0^80 w^22 s^55 i^10 k0^32
                          (q0-F)^2(s-w),
[P_product2]_198 = 9216 q0^78 w^22 s^54 i^10 k0^32(q0-F)^2,
[P_product4]_190 = 18432 q0^74 w^18 s^52 i^10 k0^32
                          (q0-F)^2(w^2-6s^2).
\end{verbatim}

They have198,99,198 terms respectively and are nonzero for each actual
positive beta. All terms, not a selected numerical ray, are
independently collected. Every gate is covered by an exact free-ring
cut, all registers, inputs and ports are live, and both alternative
operators at every gate are rejected:194,222,238 corruptions per tuple.
Normal and optimized Python run the same checks, propagate their mode to
child replays, and produce equal receipts. Relocated replay uses only
this package.
